\section{Linear Algebra Fundamentals}
\label{app:linear_algebra}

This section reviews the elements of linear algebra used
in this work, beginning with the Hermitian basis matrices.
%
The spectral decomposition is used
to derive the axis-angle representations of rotations and boosts,
%
the symmetry properties of which are used to identify
matrices that commute and the eigenmatrices of a congruence transformation.
%
Finally, the symmetry between \eqns{indexOuterBilinear}{and}{indexSpinorBilinear} is used to define a pure Mueller matrix and its associated Jones matrix, and an example of a depolarizing impure Mueller matrix is presented.

%%%%%%%%%%%%%%%%%%%%%%%%%%%%%%%%%%%%%%%%%%%%
%%%%%%%%%%%%%%%%%%%%%%%%%%%%%%%%%%%%%%%%%%%%

\subsection{Hermitian Basis Matrices}
\label{app:basis}

The Hermitian basis matrices \eqnp{Pauli} have a number of useful algebraic properties that
enable meaningful geometric interpretations of the equations used in this work.
%
When the basis matrices are indexed with a Greek character, the index is understood
to span all dimensions including the identity; e.g.,  $\mu \in \{0,1,2,3\}$.  When Latin characters are used,
the index is understood to span only the Pauli matrces; e.g.,  $m \in \{1,2,3\}$.
%
The following identities prove useful in this work.
\begin{eqnarray}
    \tr{\pauli{i}}&=&0 \label{eqn:Pauli_trace} \\
    %
    |\pauli{i}|&=&-1 \label{eqn:Pauli_determinant} \\
    %
    |\pauli{\mu} + \pauli{\nu} | &=& |\pauli{\mu}| + |\pauli{\nu}|; \ensuremath{\quad \mu \ne \nu} \label{eqn:basis_sum_determinant} \\
    %
    \pauli{\mu} \dcbo \pauli{\nu} &=& 2 \delta_{\mu\nu} \label{eqn:basis_orthogonal} \\
    %
    \pauli{i}\,\pauli{j} &=& \delta_{ij}\,\pauli{0} + i\epsilon_{ijk}\,\pauli{k} \label{eqn:Pauli_product}
\end{eqnarray}
%
In the last identity, $\epsilon_{ijk}$ is the permutation symbol and summation over the index $k$ is implied.
%
\Eqn{Pauli_trace} expresses the traceless property of the Pauli matrices.
%
\Eqns{Pauli_determinant}{and}{basis_sum_determinant} can be used to show that
\begin{equation}
| x_\mu \pauli{\mu} | = x_\mu^2 |\pauli{\mu}| = x_0^2 -x_1^2 -x_2^2 - x_3^2.
\label{eqn:determinant_is_invariant}
\end{equation}
%
That is, the determinant of a linear combination of the Hermitian
basis matrices is equal to the invariant interval of the 
four-vector of the scalar coefficients, $x_\mu$.
%
% (Note that, for any $2\times2$ matrix \JM{A} and any scalar $c$, $|cA| = c^2|A|$.)
%
\Eqn{basis_orthogonal} states that the Hermitian basis matrices
are mutually orthogonal with respect to the trace inner product.
%
The first term on the right-hand side of \eqn{Pauli_product} identifies
the Pauli matrices as the square roots of $\pauli{0}$.  
The second term describes how the Pauli matrices
behave like the basis vectors of a right-handed coordinate system, 
where $\hat{\Pv{v}}_i \times \hat{\Pv{v}}_j = \epsilon_{ijk}\,\hat{\Pv{v}}_k$.
%
Using this identity, the product of two matrices,
\begin{equation}
\begin{split}
\JM{A}\JM{B} &=
 (a\,\pauli{0}+\Pv{a}\cdot\Pv{\sigma})(b\,\pauli{0} + \Pv{b}\cdot\Pv{\sigma}) \\
        &=
 (ab + \Pv{a}\cdot\Pv{b})\pauli{0} + (a\Pv{b} + b\Pv{a}
    + \Ci \Pv{a}\times\Pv{b}) \cdot \Pv{\sigma},
    \label{eqn:multiplication_rule} 
\end{split}
\end{equation}
where $\Pv{\sigma}=\left( \pauli{1},  \pauli{2}, \pauli{3} \right)$ is
the Pauli vector that appears in \eqns{Boost}{and}{Rotation}.

\subsection{Basis Transformations}
\label{app:basis_transformations}

A pair of orthonormal receptors, $\hat{\Jrow{r}}_0$ and $\hat{\Jrow{r}}_1$, have unit length and are orthogonal; that is, they satisfy
%
\begin{equation}
\hat{\Jrow{r}}_i \cdot \hat{\Jrow{r}}_j^\dagger = \delta_{i,j}.
\label{eqn:orthonormality}
\end{equation}
%
The projection of the electric field vector onto a pair of orthonormal receptors defines a basis transformation that is equivalent to a unitary Jones matrix.
%
\begin{proof}
\label{proof:orthonormal_receptors}
%Orthonormal receptors satisfy
%$$\hat{\Jrow{r}}_i \cdot \hat{\Jrow{r}}_j^\dagger = \delta_{i,j}$$.
Let 
$$\JM{R} = \left( \begin{array}{c}
		\hat{\Jrow{r}}_0 \\
		\hat{\Jrow{r}}_1
	\end{array} \right)$$
such that 
\begin{align*}
\left\{ \JM{R} \JM{R}^\dagger \right\}_i^j 
& = \left\{ \JM{R} \right\}_i^k \left\{ \JM{R}^\dagger \right\}_k^j
  & \peqn{matrix_product} \\
& = \left\{ \hat{\Jrow{r}}_i \right\}^k \left\{ \hat{\Jrow{r}}_j^\dagger \right\}_k  \\
& = \hat{\Jrow{r}}_i \cdot \hat{\Jrow{r}}_j^\dagger 
  & \peqn{scalar_product} \\
& = \delta_{i,j} 
  & \peqn{orthonormality}
\end{align*}
That is, $\JM{R} \JM{R}^\dagger = \pauli{0}$; 
% is equal to the identity matrix.
therefore, $\JM{R}^\dagger=\JM{R}^{-1}$.
\end{proof}
%
In general, a unitary matrix has a determinant equal to $\exp(\Ci\phi)$, and
multiplication by $\exp(-\Ci\phi/2)$ yields a unimodular matrix
(with $|\JM{R}|=1$).
%
Basis transformations defined by unitary matrices are used extensively throughout both \Sec{signal_path} and the following sections of the appendix.  

\subsection{Spectral Decomposition}
\label{app:spectral_decomposition}

\iffalse
\Eqns{Boost}{and}{Rotation} can be derived using the spectral decomposition,
% of a linear combination of the Pauli matrices,
which is a special case of the eigendecomposition.
\fi
%
The eigenvalues $\lambda$ and eigenvectors $\Jcol{v}$ of
a matrix \JM{A} satisfy the following (logically-equivalent) equations,
%
\begin{equation}
    \label{eqn:eigen}
    \JM{A}\Jcol{v} = \lambda\Jcol{v} \quad \Leftrightarrow \quad
    \left(\JM{A} - \lambda\JM{I} \right) \Jcol{v} = 0,
\end{equation}
where \JM{I} is the identity matrix with the dimensions of \JM{A}.
%
\Eqn{eigen} can be expressed simultaneously for all eigenvalue/eigenvector pairs by $\JM{A} \JM{R} = \JM{R} \cohM{\Lambda}$, where the $k^\mathrm{th}$
column of {\bf{R}} is the unit eigenvector $\hat{\Jcol{v}}_k$, and the diagonal matrix $\cohM{\Lambda}$ is defined by $\Lambda_k^j = \delta_{j,k} \lambda_k$.
%
This can be rearranged to yield the eigendecomposition of \JM{A},
%
\begin{equation}
\JM{A} = \JM{R} \cohM{\Lambda} \JM{R}^{-1} \label{eqn:eigendecomposition}.
\end{equation}

If \JM{A} is normal (i.e., $\JM{A}^\dagger\JM{A}=\JM{A}\JM{A}^\dagger$),
and $\Jcol{v}$ is an eigenvector of \JM{A} with associated eigenvalue $\lambda$, 
then $\Jcol{v}$ is also an eigenvector of $\JM{A}^\dagger$ with associated eigenvalue $\lambda^*$.

\begin{proof}
\label{proof:normal_eigenvalue}
Define $\JM{L}=\JM{A} - \lambda\JM{I}$ and consider
\begin{align*}
|\JM{L}^\dagger \Jcol{v}|^2
  &= \left( \JM{L}^\dagger \Jcol{v} \right)^\dagger \left( \JM{L}^\dagger \Jcol{v} \right) \\
  &= \Jcol{v}^\dagger \JM{L} \JM{L}^\dagger \Jcol{v} \\
  &= \Jcol{v}^\dagger \JM{L}^\dagger \JM{L} \Jcol{v} & \text{\JM{L} is normal} \\
  &= \left( \JM{L}\Jcol{v} \right)^\dagger \JM{L} \Jcol{v} \\
  &= |\JM{L} \Jcol{v}|^2 \\
  &= |\left(\JM{A} - \lambda\JM{I} \right) \Jcol{v}|^2 = 0 & \peqn{eigen}
\end{align*}
Therefore $\JM{L}^\dagger \Jcol{v} = \left(\JM{A}^\dagger - \lambda^*\JM{I} \right) \Jcol{v} = 0.$
\end{proof}
\noindent
Furthermore, if \JM{A} is normal, then \eqn{eigendecomposition} is equivalent to a
congruence transformation by a unitary matrix known as its spectral decomposition,
%
\begin{equation}
\label{eqn:congruence_eigendecomposition}
\JM{A} = \JM{R} \cohM{\Lambda} \JM{R}^\dagger.
\end{equation}
%

\begin{proof}
\label{proof:spectral_decomposition}
If \JM{A} is normal, then
\begin{align*}
( \JM{A}^\dagger \hat{\Jcol{v}}_i )^\dagger \hat{\Jcol{v}}_j
&= \hat{\Jcol{v}}_i^\dagger \JM{A} \hat{\Jcol{v}}_j  \\
%
%
( \lambda_i^* \hat{\Jcol{v}}_i)^\dagger \cdot \hat{\Jcol{v}}_j
&= \hat{\Jcol{v}}_i^\dagger \cdot ( \lambda_j \hat{\Jcol{v}}_j)
 & \text{\Proof{normal_eigenvalue}} \\
%
%
\lambda_i \hat{\Jcol{v}}_i^\dagger \cdot \hat{\Jcol{v}}_j
&= \lambda_j \hat{\Jcol{v}}_i^\dagger \cdot \hat{\Jcol{v}}_j \\
%
%
(\lambda_i-\lambda_j) \hat{\Jcol{v}}_i^\dagger \cdot \hat{\Jcol{v}}_j &= 0
\end{align*}
If $\lambda_i \ne \lambda_j$, then
$\hat{\Jcol{v}}_i^\dagger \cdot \hat{\Jcol{v}}_j = 0$ and
\begin{align*}
\left\{ \JM{R}^\dagger \JM{R} \right\}_i^j 
& = \left\{ \JM{R}^\dagger \right\}_i^k \left\{ \JM{R} \right\}_k^j \\
& = \left\{ \hat{\Jcol{v}}_i^\dagger \right\}^k \left\{ \hat{\Jcol{v}}_j \right\}_k 
  = \hat{\Jcol{v}}_i^\dagger \cdot \hat{\Jcol{v}}_j = \delta_{i,j}
\end{align*}
That is, $\JM{R}^\dagger \JM{R} = \JM{I}$; therefore, $\JM{R}^\dagger=\JM{R}^{-1}.$
\end{proof}

In the natural basis defined by ${\bf{R}}^\dagger$, the matrix $\JM{A}$ becomes the diagonal matrix,
%
\begin{equation}
\label{eqn:diagonal_decomposition}
\cohM{\Lambda} = \JM{R}^\dagger \JM{A} \JM{R}.
\end{equation}
%
Both Hermitian and unitary matrices are normal, and each takes
diagonal form in the natural basis defined by its eigenvectors.  Consequently, any Hermitian matrix can be diagonalized such that it is equivalent to a differential gain transformation (see \Sec{differential_gain}) in its natural basis.
%
Similarly, any unitary matrix is equivalent to differential phase (see \Sec{differential_phase}) in its natural basis.
% (e.g., see \Proof{Faraday}).
%
This is applied in \Proof{Faraday_rotation}, where rotation about the line of sight is shown to be equivalent to differential phase between the circularly-polarized natural modes of a magnetized cold plasma.



The spectral decomposition can also be written as a linear combination of outer products,
\begin{equation}
\JM{A} = \lambda_k\, \hat{\Jcol{v}}_k \otimes \hat{\Jcol{v}}_k^\dagger.
\label{eqn:spectral_decomposition}
\end{equation}

\begin{proof}

\begin{align*}
A_i^j
&= \left\{ {\bf{R}} \cohM{\Lambda} {\bf{R}}^\dagger \right\}_i^j 
    & \peqn{congruence_eigendecomposition} \\
&=  \left\{ {\bf{R}} \right\}_i^k
    \left\{ \cohM{\Lambda} \right\}_k^l
    \left\{ {\bf{R}}^\dagger \right\}_l^j 
    & \peqn{matrix_product} \\
&= \left\{ \hat{\Jcol{e}}_k \right\}_i \; \delta_{k,l} \; \lambda_k \left\{ \hat{\Jcol{e}}_l^\dagger \right\}^j \\
&= \lambda_k \;\left\{ \hat{\Jcol{e}}_k \right\}_i \left\{ \hat{\Jcol{e}}_k^\dagger \right\}^j \\
&= \lambda_k \; \left\{ \hat{\Jcol{e}}_k \otimes \hat{\Jcol{e}}_k^\dagger \right\}_i^j
    & \peqn{implicit_tensor_product}
\label{eqn:eigen_sum}
\end{align*}
%
\end{proof}
\noindent
If the eigenvalues are distinct, then
\begin{equation}
\label{eqn:projection}
\JM{P}_k=\hat{\Jcol{v}}_k \otimes \hat{\Jcol{v}}_k^\dagger
\end{equation}
form a complete set of mutually orthogonal projection matrices; i.e.,
\begin{equation}
\label{eqn:idempotent_orthogonal}
\JM{P}_i \JM{P}_j = \delta_{i,j} \JM{P}_j
\end{equation}
and
\begin{equation}
\sum_i \JM{P}_i = \JM{I}
\label{eqn:projection_complete}
\end{equation}
%
\begin{proof}
\label{proof:orthogonal_idempotent}
$\JM{P}_k$ are orthogonal and idempotent.
\begin{align*}
\JM{P}_i \JM{P}_j &= 
(\hat{\Jcol{v}}_i  \hat{\Jcol{v}}_i^\dagger)
(\hat{\Jcol{v}}_j  \hat{\Jcol{v}}_j^\dagger) 
    & \peqn{implicit_tensor_product} \\
&= \hat{\Jcol{v}}_i ( \hat{\Jcol{v}}_i^\dagger
\hat{\Jcol{v}}_j ) \hat{\Jcol{v}}_j^\dagger
    & \text{Associativity} \\
&= \delta_{i,j} \hat{\Jcol{v}}_i \hat{\Jcol{v}}_j^\dagger
   & \hat{\Jcol{v}}_i^\dagger \hat{\Jcol{v}}_j = \delta_{i,j} \\
&= \delta_{i,j} \JM{P}_j & \peqn{projection}
\end{align*}
\end{proof}
%
\begin{proof}
$\JM{P}_k$ form a complete set.

Let $\Jcol{w} = c_j \hat{\Jcol{v}}_j$ and consider
\begin{align*}
\sum_i \JM{P}_i \Jcol{w} 
&= \left( \hat{\Jcol{v}}_i \hat{\Jcol{v}}_i^\dagger \right)
\left( c_j \hat{\Jcol{v}}_j \right) 
  & \peqn{projection} \\
&= c_j \hat{\Jcol{v}}_i \hat{\Jcol{v}}_i^\dagger \hat{\Jcol{v}}_j 
  & \text{Commutativity} \\
&= c_j \hat{\Jcol{v}}_i \delta_{i,j} 
  & \hat{\Jcol{v}}_i^\dagger \hat{\Jcol{v}}_j = \delta_{i,j} \\
& = c_j \hat{\Jcol{v}}_j \\
& = \Jcol{w}
\end{align*}
\end{proof}



The $2 \times 2$ Hermitian coherency matrix has two eigenvalues, $\lambda_k = (S_0 \pm |\Pv{S}|)/2$, and
can be expressed as a spectral decomposition,
\begin{equation}
\begin{split}
\label{eqn:spectral_decomposition_of_coherency_matrix}
\cohM{\rho} 
&=  \lambda_0\, \hat{\Jcol{e}}_0 \otimes \hat{\Jcol{e}}_0^\dagger
    + \lambda_1\, \hat{\Jcol{e}}_1 \otimes \hat{\Jcol{e}}_1^\dagger \\
&=  \lambda_0\,\tilde{\cohM{\rho}}_0 + \lambda_1\,\tilde{\cohM{\rho}}_1,
\end{split}
\end{equation}
where $\tilde{\cohM{\rho}}_k = \hat{\Jcol{e}}_k \otimes \hat{\Jcol{e}}_k^\dagger$ correspond to purely polarized states,
as in \eqn{instantaneous_coherency}.
%
That is, any partially polarized state can be represented as an incoherent superposition of a pair of purely polarized states
that are orthogonally polarized ($\hat{\Jcol{e}}_0^\dagger\hat{\Jcol{e}}_1 = 0$ and $\tilde{\cohM{\rho}}_0 \tilde{\cohM{\rho}}_1 = 0$).
%
If the signal is unpolarized, then the eigenvalues are equal and
any non-zero vector is an eigenvector.
%
If $\cohM{\rho}$ is purely polarized, then one of the eigenvalues equals zero, and the polarization state is completely described by the Jones vector associated with the non-zero eigenvalue.

In the natural basis defined by ${\bf{R}}^\dagger$, the eigenvalues
$\lambda_m$ are equal to the variances of two uncorrelated signals
received by orthogonally polarized receptors described by
the Hermitian transposes of the eigenvectors.  In this basis,
%
the total intensity, $S_0=\lambda_0+\lambda_1$; the
polarized intensity, $S_1=|\Pv{S}|=\lambda_0-\lambda_1$; and
$S_2=S_3=0$.
%
That is, ${\bf{R}}^\dagger$ rotates the basis such that the
polarization vector points along $S_1$.  In this basis, it is clear that the degree of polarization
%
\begin{equation}
\label{eqn:degree_of_polarization_eigen}
p = \frac{|\Pv{S}|}{S_0} = \frac{|\lambda_0-\lambda_1|}{\lambda_0+\lambda_1},
\end{equation}
%
is equal to zero when the intensities of the orthogonal modes are equal ($\lambda_0 = \lambda_1$);
in contrast, $p=1$ when the signal consists of a single purely polarized mode (one of the eigenvalues is zero).

\subsection{Axis-angle Representation}
\label{app:fundamental_transformations}

\Eqns{Boost}{and}{Rotation} are defined using the matrix exponential, which is given by the power series,
%
\begin{equation}
    \exp(\JM{A}) \equiv \sum_{k=0}^\infty \frac{\JM{A}^k}{k!}.
    \label{eqn:matrix_exponential}
\end{equation}
%
Here, $\JM{A}^k$ is the $k^\mathrm{th}$ power of $\JM{A}$; 
$\JM{A}^0 = \JM{I}$, where \JM{I} is the identity matrix with the dimensions of \JM{A};
and $k!$ is the factorial of $k$.
%
%
If \JM{A} is normal, then its spectral decomposition
% the powers of \JM{A} can be expressed using its spectral decomposition 
\eqnsp{spectral_decomposition}{and}{projection},
\begin{equation}
    \JM{A} = \lambda_k \JM{P}_k,
\end{equation}
such that
\begin{equation}
    \JM{A}^n = (\lambda_k)^n \JM{P}_k.
    \label{eqn:matrix_power}
\end{equation}

\begin{proof}
Using induction, start with the base case and apply \eqn{projection_complete} to yield
\begin{align*}
\JM{A}^0 = (\lambda_k)^0 \JM{P}_k = \sum_k \JM{P}_k = \JM{I}
\end{align*}
If $\JM{A}^n = (\lambda_k)^n \JM{P}_k$, then
\begin{align*}
\JM{A}^{n+1} &= \JM{A}\JM{A}^n \\
&= \left( \lambda_j \JM{P}_j \right) \left( (\lambda_k)^n \JM{P}_k \right)  \\
&= \lambda_j (\lambda_k)^n \JM{P}_j \JM{P}_k   \\
&= \lambda_j (\lambda_k)^n \delta_{jk} \JM{P}_k & \peqn{idempotent_orthogonal} \\
&= (\lambda_k)^{n+1} \JM{P}_k. & \hfill \square
\end{align*}
%
\end{proof}

\noindent
Substitute \eqn{matrix_power} into \eqn{matrix_exponential} to yield
\begin{align}
\exp(\JM{A}) 
% & = \sum_{k=0}^\infty \frac{1}{k!} \left( (\lambda_j)^n \JM{P}_j \right) \\
& = \left( \sum_{k=0}^\infty \frac{(\lambda_j)^n}{k!} \right) \JM{P}_j 
  = \exp(\lambda_j) \JM{P}_j.
\end{align}

To arrive at \eqns{Boost}{and}{Rotation}, consider the eigenvalues of 
$\JM{A} = \Pv{a}\cdot\Pv{\sigma}$, which must satisfy
%
% \eqnp[via ]{determinant_is_invariant}
\begin{equation}
    |\JM{A} - \lambda \JM{I}| = \left|\Pv{a}\cdot\Pv{\sigma} - \lambda \pauli{0} \right| = \lambda^2 - |\Pv{a}|^2 = 0.
\end{equation}
That is, there are two eigenvalues given by $\lambda = \pm |\Pv{a}|$.
%
Let $a=|\Pv{a}|$ and $\hat{\Pv{a}}=\Pv{a}/a$, and use the spectral decomposition,
\begin{equation}
\JM{A} = a\JM{P}_0 -a\JM{P}_1,
\end{equation}
to show that
\begin{equation}
\JM{P}_0 + \JM{P}_1 = \pauli{0} \quad \text{and} \quad
\JM{P}_0 - \JM{P}_1 = \hat{\Pv{a}}\cdot\Pv{\sigma};
\end{equation}
therefore
\begin{align}
e^{\Pvs{a}\cdot\Pvs{\sigma}}
&= \exp(a) \JM{P}_0 + \exp(-a) \JM{P}_1 \nonumber \\
&= \cosh(a) (\JM{P}_0 + \JM{P}_1) 
+ \sinh(a)(\JM{P}_0 - \JM{P}_1) \label{eqn:Pauli_exponential} \\
&= \cosh(a) \pauli{0}
+ \sinh(a) \hat{\Pv{a}}\cdot\Pv{\sigma}. \nonumber
\end{align}
%
Setting $\Pv{a} = \beta \hat{\Pv{m}}$ in \eqn{Pauli_exponential} 
yields the axis-angle representation of Hermitian matrices \eqnp{Boost};
likewise, setting $\Pv{a} = \Ci \phi \hat{\Pv{n}}$ yields the axis-angle representation of unitary matrices \eqnp{Rotation}.


%%%%%%%%%%%%%%%%%%%%%%%%%%%%%%%%%%%%%%%%%%%%
%%%%%%%%%%%%%%%%%%%%%%%%%%%%%%%%%%%%%%%%%%%%

\subsection{Congruence Eigenmatrices}
\label{app:symmetry}

The unit vector in the axis-angle representation of a matrix defines an axis of symmetry in the three-dimensional space of the Stokes polarization vector.
%
This symmetry axis can be exploited to rearrange and simplify matrix equations and identify degenerate systems of equations \citep[as in Appendix B of][and \App{degeneracy} of this paper]{van04c}.
%
For example, two matrices 
% \JM{A} and \JM{B} 
commute when their symmetry axes are collinear (parallel or anti-parallel).
%
\begin{proof}
Given $$\JM{A} = a\,\pauli{0} + \Pv{a} \cdot \Pv{\sigma} \quad\text{and}\quad
\JM{B} = b\,\pauli{0} + \Pv{b} \cdot \Pv{\sigma},$$
 \eqn{multiplication_rule} shows that $\JM{A}\JM{B} = \JM{B}\JM{A}$ if and only if $\Pv{a} \times \Pv{b} = 0$.
\end{proof}
%
\iffalse
\noindent
Furthermore, the axis of symmetry of the product of two matrices that commute
is collinear with the symmetry axes of both matrix factors;
e.g.,  
\begin{equation}
    (\JM{A}\JM{B})\JM{A} = \JM{A}(\JM{A}\JM{B}) \quad \text{and} \quad (\JM{A}\JM{B})\JM{B} = \JM{B}(\JM{A}\JM{B}).
    \label{eqn:collinear_product}
\end{equation}
\fi
Furthermore, if \JM{A} and \JM{B} commute, then they have common eigenvectors.
\begin{proof}
\label{proof:common_eigenvalue}
If $\JM{A}\JM{B} = \JM{B}\JM{A}$ and $\JM{A}\Jcol{v} = \lambda \Jcol{v}$,
then
\begin{align*}
\JM{A}\JM{B}\Jcol{v} &= \JM{B}\JM{A}\Jcol{v} = \lambda \JM{B}\Jcol{v}.
\end{align*}
That is, $\JM{B}\Jcol{v}$ is an eigenvector of \JM{A} with the same eigenvalue $\lambda$;
therefore, $\JM{B}\Jcol{v}=\mu\Jcol{v}$.
\end{proof}


The coherency matrix has an axis of symmetry defined by its associated Stokes polarization vector, which facilitates the identification of the eigenmatrices of a congruence transformation.
%
A matrix $\cohM{\rho}$ is a congruence eigenmatrix of an invertible matrix $\JM{J}$ with real-valued congruence eigenvalue $\kappa$ if it satisfies the relation
%
\begin{equation}
\JM{J}\cohM{\rho}\JM{J}^\dagger = \kappa \cohM{\rho}.
\end{equation}
%
The eigenvectors of a Jones matrix define the singular congruence eigenmatrices of that Jones matrix.
%
\begin{proof}
\label{proof:eigenvector_to_eigenmatrix}
If $\JM{J} \Jcol{e} = \lambda \Jcol{e}$, where $\lambda \in \mathbb{C}$, then
\begin{align*}
\JM{J}\tilde{\cohM{\rho}}\,\JM{J}^\dagger
&= \JM{J} \Jcol{e} \otimes \Jcol{e}^\dagger \JM{J}^\dagger
    & \peqn{instantaneous_coherency} \\
&= \JM{J} \Jcol{e} \otimes  (\JM{J}\Jcol{e})^\dagger 
    & (\JM{A}\JM{B})^\dagger = \JM{B}^\dagger \JM{A}^\dagger \\
&= \lambda \, \Jcol{e} \otimes \lambda^* \Jcol{e}^\dagger \\
&= |\lambda|^2 \tilde{\cohM{\rho}} & \tilde{\cohM{\rho}} \equiv \Jcol{e} \otimes \Jcol{e}^\dagger \\
&= \kappa \tilde{\cohM{\rho}} & \kappa \equiv |\lambda|^2 \in \mathbb{R}
\end{align*}
\end{proof}
%
\noindent
Accordingly, a Jones matrix has two congruence eigenmatrices.\footnote{
%
Although the Mueller matrix associated with \JM{J}
has four eigenvectors, only two of these
describe physical polarization states.
%
For example, two of the eigenvectors of a rotation are complex-valued.  
%
Two of the eigenvectors of a Lorentz boost have polarization vectors that lie 
in the plane that is normal to the boost axis; however, these have $S_0 = 0$, which is not physical.   }
%
% e.g.,  https://physics.stackexchange.com/questions/108520/what-are-the-eigenvalues-of-the-lorentz-matrix
%
% Is it possible to prove this?  For example, by treating linear operators as rank 4 tensors and showing that only 2 of the 4 eigenmatrices can have positive determinant?
%
% The two \emph{physical} eigenmatrices of a congruence transformation correspond to the purely polarized eigenvectors of the Jones matrix.
%
Conversely, if a singular coherency matrix
%
$\hat{\cohM{\rho}}=\Jcol{e} \otimes \Jcol{e}^\dagger$ 
%
is a congruence eigenmatrix of \JM{J}, then $\Jcol{e}$
is an eigenvector of \JM{J}.
%
\begin{proof}
%
\label{proof:eigenmatrix_to_eigenvector}
%
If $\hat{\cohM{\rho}}=\Jcol{e} \otimes \Jcol{e}^\dagger$ and $\JM{J}\hat{\cohM{\rho}}\JM{J}^\dagger = \kappa \hat{\cohM{\rho}}$, then
%
$$ \JM{J} \left( \Jcol{e} \otimes \Jcol{e}^\dagger \right) \JM{J}^\dagger
=
\JM{J} \Jcol{e} \otimes  (\JM{J}\Jcol{e})^\dagger 
=
\kappa \Jcol{e} \otimes \Jcol{e}^\dagger, $$
and $\JM{J} \Jcol{e}$ and $\Jcol{e}$ must be collinear; i.e.,
$\JM{J} \Jcol{e} = \lambda \Jcol{e}$, where $|\lambda|^2 = \kappa$.
\end{proof}
%
\noindent
Furthermore, if
%
$\hat{\cohM{\rho}}=\Jcol{e} \otimes \Jcol{e}^\dagger$ 
%
commutes with \JM{J}, then $\hat{\cohM{\rho}}$
is a congruence eigenmatrix of $\JM{J}$.
%
\begin{proof}
\label{proof:if_commutes_then_eigen}
If $\JM{J} \tilde{\cohM{\rho}} = \tilde{\cohM{\rho}} \JM{J}$, then 
$\tilde{\cohM{\rho}}$ and $\JM{J}$ have a common eigenvector (\Proof{common_eigenvalue}).
Therefore,
$$\JM{J} \Jcol{e} = \lambda \, \Jcol{e}$$
%
where $\Jcol{e}$ is the only eigenvector of $\tilde{\cohM{\rho}}=\Jcol{e} \otimes \Jcol{e}^\dagger$, and
%
(via \Proof{eigenvector_to_eigenmatrix})
%
$$\JM{J} \tilde{\cohM{\rho}}\, \JM{J}^\dagger = \kappa\tilde{\cohM{\rho}},$$
where $\kappa=|\lambda|^2$.
\end{proof}
%

\noindent
Finally, if 
$\JM{J}$ is a normal matrix
and 
$\hat{\cohM{\rho}}$ is a congruence eigenmatrix of $\JM{J}$,
then 
$\hat{\cohM{\rho}}$ commutes with $\JM{J}$.

\begin{proof}
\label{proof:if_eigen_then_commutes}
If $\JM{J}$ is normal,
$\hat{\cohM{\rho}}=\Jcol{e} \otimes \Jcol{e}^\dagger$ 
and
$\JM{J} \tilde{\cohM{\rho}}\, \JM{J}^\dagger = \kappa\tilde{\cohM{\rho}},$
then 
$\JM{J} \Jcol{e} = \lambda \Jcol{e}$ (\Proof{eigenmatrix_to_eigenvector})
and
$$\JM{J} \tilde{\cohM{\rho}} 
  = \JM{J} \Jcol{e} \otimes \Jcol{e}^\dagger 
  = \lambda \, \Jcol{e} \otimes \Jcol{e}^\dagger 
  = \lambda \tilde{\cohM{\rho}}. $$

Also
\begin{align*}
 \tilde{\cohM{\rho}} \JM{J}
  &=  \Jcol{e} \otimes \Jcol{e}^\dagger \JM{J} \\
  &= \Jcol{e} \otimes \left( \JM{J}^\dagger \Jcol{e} \right)^\dagger  \\
  &= \Jcol{e} \otimes \left( \lambda^* \Jcol{e} \right)^\dagger & \text{\Proof{normal_eigenvalue}} \\
  &= \lambda \, \Jcol{e} \otimes  \Jcol{e}^\dagger \\
  &= \lambda \tilde{\cohM{\rho}}
\end{align*}
Therefore, $\JM{J} \tilde{\cohM{\rho}} = \tilde{\cohM{\rho}} \JM{J}.$
\end{proof}

\noindent
%
In conclusion, if $\JM{J}$ is normal,  then $\hat{\cohM{\rho}}$ is a congruence eigenmatrix of $\JM{J}$ if and only if $\hat{\cohM{\rho}}$ commutes with $\JM{J}$.
%
Both unitary and Hermitian matrices are normal and their congruence eigenmatrices are explored in the following sections.

%%%%%%%%%%%%%%%%%%%%%%%%%%%%%%%%%%%%%%%%%%%%
%%%%%%%%%%%%%%%%%%%%%%%%%%%%%%%%%%%%%%%%%%%%
\subsubsection{ Unitary Matrices}
\label{app:eigen_unitary}

If $\cohM{\rho}$ commutes with \rotat, then
\begin{equation}
\rotat \, \cohM{\rho} \, \vRotation^\dagger(\phi)
= \cohM{\rho}  \rotat \vRotation^\dagger(\phi) 
= \cohM{\rho};
\end{equation}
that is, $\cohM{\rho}$ is a congruence eigenmatrix of \rotat\ with associated congruence eigenvalue equal to unity.
% 
Conversely, if $\cohM{\rho}$ is a congruence eigenmatrix of \rotat, then
\begin{equation}
\begin{split}
\rotat \, \cohM{\rho} \, \vRotation^\dagger(\phi) & = \kappa \cohM{\rho}  \\
\rotat \, \cohM{\rho} \, \vRotation^\dagger(\phi)\, \rotat & = \kappa \cohM{\rho} \, \rotat \\
\rotat \, \cohM{\rho} & = \kappa \cohM{\rho} \, \rotat
\end{split}
\end{equation}
and $\cohM{\rho}$ commutes with \rotat\ if $\kappa = 1$.
\noindent
The congruence eigenmatrices of a unitary matrix share a single 
degenerate congruence eigenvalue, $\kappa=1$; therefore, any linear combination of the congruence eigenmatrices
is also a congruence eigenmatrix of the transformation.
%
Considering the equivalent three-dimensional rotation of the Stokes polarization vector $\Pv{S}$;
any vector that lies along the axis of rotation remains unchanged by that rotation, including $\Pv{S}=0$ (unpolarized radiation).
%


%%%%%%%%%%%%%%%%%%%%%%%%%%%%%%%%%%%%%%%%%%%%
%%%%%%%%%%%%%%%%%%%%%%%%%%%%%%%%%%%%%%%%%%%%
\subsubsection{ Hermitian Matrices }
\label{app:eigen_self_adjoint}

Define $b=\cosh\beta$ and $\Pv{b}=\hat{\Pv{m}}\sinh\beta$ such that 
\begin{equation}
    \boost = b\,\pauli{0} + \Pv{b}\cdot\Pv{\sigma}.
\end{equation}
%
Furthermore, define 
$\tilde{\cohM{\rho}} = \left(\tilde{s}_0\,\pauli{0} + \tilde{\Pv{s}}\cdot \Pv{\sigma}\right)/2$,
and note that $|\tilde{\Pv{s}}| = \tilde{s}_0$ because $\tilde{\cohM{\rho}}$ is singular / purely polarized. 
%
If $\tilde{\cohM{\rho}}$ and \boost\ commute, then their axes of
symmetry are collinear; i.e.,
\begin{equation}
\begin{split}
\hat{\Pv{m}} \times \tilde{\Pv{s}} &= 0, \\
\hat{\Pv{m}} \cdot \tilde{\Pv{s}} &= \pm|\tilde{\Pv{s}}| = \pm\tilde{s}_0, \\
\hat{\Pv{m}} &= \pm \, \tilde{\Pv{s}} / \tilde{s}_0.
\end{split}
\end{equation}
%
Using \eqn{multiplication_rule},
\begin{equation} \begin{split}
\boost \tilde{\cohM{\rho}} 
&= (b \tilde{s}_0 + \Pv{b} \cdot \tilde{\Pv{s}})\pauli{0}/2 + (b\tilde{\Pv{s}} + \tilde{s}_0\Pv{b})\cdot\Pv{\sigma}/2 \\
&= (b \tilde{s}_0 \pm |\Pv{b}| \tilde{s}_0)\pauli{0}/2 + (b\tilde{\Pv{s}} \pm |\Pv{b}| \tilde{\Pv{s}})\cdot\Pv{\sigma}/2 \\
&= (b \pm |\Pv{b}|) \tilde{s}_0  \pauli{0}/2 + (b \pm |\Pv{b}|) \tilde{\Pv{s}}\cdot\Pv{\sigma}/2 \\
&= (b \pm |\Pv{b}|) \tilde{\cohM{\rho}} \\
&= (\cosh\beta \pm \sinh\beta) \tilde{\cohM{\rho}} \\
&= e^{\pm\beta} \tilde{\cohM{\rho}}.
\end{split} \end{equation}
Therefore,
\begin{equation} \begin{split}
\boost \, \tilde{\cohM{\rho}} \, \vBoost^\dagger(\beta)
&= \boost \, \tilde{\cohM{\rho}} \, \boost \\
&= \vBoost^2(\beta) \tilde{\cohM{\rho}} \\
&= e^{\pm2\beta} \tilde{\cohM{\rho}}.
\end{split} \end{equation}
%
Hermitian matrices have two distinct congruence eigenvalues; 
$\kappa=e^{2\beta}$ when the polarization vector $\tilde{\Pv{s}}$
is parallel to the boost axis $\hat{\Pv{m}}$,
and
$\kappa=e^{-2\beta}$ when $\tilde{\Pv{s}}$
is anti-parallel to $\hat{\Pv{m}}$.
%
Therefore, only these two purely polarized states are congruence eigenmatrices of the transformation.

%%%%%%%%%%%%%%%%%%%%%%%%%%%%%%%%%%%%%%%%%%%%
%%%%%%%%%%%%%%%%%%%%%%%%%%%%%%%%%%%%%%%%%%%%

\subsection{ Pure Mueller Matrices }
\label{app:pure_Mueller_matrices}
%
Various authors \cite[e.g.,][]{barakat81,sim82,clo86} have considered the necessary and sufficient conditions that must be satisfied by the elements of a Mueller matrix for it to have an equivalent Jones matrix representation.
%
\cite{sim82} and \cite{clo86} present the most intuitive geometric constraints on 
a pure Mueller matrix based on its equivalent target coherency matrix \citep{clo86}.

For a Jones matrix \JM{J}, the equivalent $4 \times 4$ Hermitian target coherency matrix $\MM{N}$ is defined via the rank 4 tensor $\JT{N}=\outerBilinear{\bf J}{\bf J}^\dagger$, such that
\eqn{tensor_to_Mueller} yields
%
\begin{equation}
\begin{split}
\label{eqn:target_coherency_matrix}
N_\irow^\icol
& = \frac{1}{2} \dc{\pauli{\irow}}{\dc{\left( \outerBilinear{\bf J}{\bf J}^\dagger \right)}{\pauli{\icol}}} \\
& = \frac{1}{2} \dc{\pauli{\irow}}{{\bf J}} \left( \dc{\pauli{\icol}}{{\bf J}^\dagger} \right) \\
& = \frac{1}{2} k_\irow k^*_\icol.
\end{split}
\end{equation}
%
Here, $k_\irow = \dc{\pauli{\irow}}{{\bf J}}$ are the components of a complex-valued target four-vector $\Sv{k}$ \citep{clo86}.
%
\Eqn{target_coherency_matrix} is equivalent to $\MM{N} = \Sv{k} \otimes \Sv{k}^\dagger / 2$ and, as
noted by \cite{sim82}, \MM{N} is a scalar multiple of a projection matrix, such that
\begin{equation}
\MM{N}^2 = \tr{\MM{N}} \MM{N}.
\label{eqn:pure_target_coherency}
\end{equation}

\begin{proof}
Define $\hat{\Sv{k}}=\Sv{k}/|\Sv{k}|$ and the projection $\MM{P} = \hat{\Sv{k}} \otimes \hat{\Sv{k}}^\dagger$ \eqnp{projection}, 
such that
\begin{equation*}
\MM{N}
 = \frac{1}{2}\Sv{k} \otimes \Sv{k}^\dagger \\
 = \frac{1}{2}|\Sv{k}|^2 \hat{\Sv{k}} \otimes \hat{\Sv{k}}^\dagger
 = \frac{1}{2}|\Sv{k}|^2 \MM{P}.
\end{equation*}
Note that 
\begin{equation*}
|\Sv{k}|^2 = k_\mu k_\mu^* = 2 \tr{\MM{N}}.
\end{equation*}
Therefore, $\MM{N} = \tr{\MM{N}} \MM{P}$ and \eqnp[via ]{idempotent_orthogonal}
\begin{equation*}
\MM{N}^2 = \tr{\MM{N}}^2 \MM{P}^2 = \tr{\MM{N}}^2 \MM{P}= \tr{\MM{N}} \MM{N}.
\end{equation*}
\end{proof}

The mapping between any Mueller matrix $\MM{M}$ and its associated target coherency matrix $\MM{N}$ exploits the symmetry of the transpose used to define the $\tilde\otimes$ operator \eqnp{indexSpinorBilinear}.
%
Let $\mathcal{T}$ represent the transpose of covariant tensor indeces, such that
$\spinorBilinear{\bf A}{\bf B} = \mathcal{T}(\outerBilinear{\bf A}{\bf B})$ and, by symmetry,
$\outerBilinear{\bf A}{\bf B} = \mathcal{T}(\spinorBilinear{\bf A}{\bf B})$.  Using this operator, the rank 4 tensors associated with \MM{M} and \MM{N} are related by $\JT{U} = \mathcal{T}(\JT{N})$, where $\JT{U}$ is defined by \eqn{Mueller_to_tensor}
and 
%
\begin{equation}
\begin{split}
\JT{N} & = \mathcal{T}(\JT{U}) \\
 & = \mathcal{T}\left( \frac{1}{2} M_\irow^\icol \outerBilinear{\pauli{\irow}}{\pauli{\icol}} \right) \\
 & = \frac{1}{2} M_\irow^\icol \spinorBilinear{\pauli{\irow}}{\pauli{\icol}},
\end{split}
\end{equation}
such that
\begin{equation}
\begin{split}
\label{eqn:Mueller_to_target_coherency}
N_\jrow^\jcol & = \frac{1}{2} \dc{\pauli{\jrow}}{\dc{ \JT{N} }{\pauli{\jcol}}} \\
& = \frac{1}{4} M_\irow^\icol \, \dc{\pauli{\jrow}}{\dc{\left( \spinorBilinear{\pauli{\irow}}{\pauli{\icol}} \right)}{\pauli{\jcol}}} \\
& = \frac{1}{4} M_\irow^\icol \, \dc{\pauli{\jrow}}{ \pauli{\irow}}{\pauli{\jcol}} {\pauli{\icol}} \\
& = \frac{1}{4} M_\irow^\icol \, \tr{ \pauli{\jrow} \pauli{\irow} \pauli{\jcol} \pauli{\icol} }. \\
\end{split}
\end{equation}
\Eqn{Mueller_to_target_coherency} can be written
as $\MM{N} = \MM{\JT{T}} \dcbo \MM{M}$, where $\MM{\JT{T}}$ is the
four-dimensional rank 4 tensor defined by
%
\begin{equation}
T_{\jrow\icol}^{\jcol\irow} = \frac{1}{4}  \tr{ \pauli{\jrow} \pauli{\irow} \pauli{\jcol} \pauli{\icol} }.
\end{equation}
By symmetry, $\MM{M} = \MM{\JT{T}} \dcbo \MM{N}$; that is, $\MM{\JT{T}}$ is an involution.
%
As for the target coherency matrix derived from a Jones matrix,
$\MM{N}$ is Hermitian.

\begin{proof}
\begin{align*}
&(N_\jrow^{\jcol})^* = (T_{\jrow\icol}^{\jcol\irow})^* (M_\irow^\icol)^* \\
&= \frac{1}{4} M_\irow^\icol \, \tr{ \pauli{\jrow} \pauli{\irow} \pauli{\jcol} \pauli{\icol} }^* & M_\irow^\icol\in\mathbb{R} \\
&= \frac{1}{4} M_\irow^\icol \, \tr{ \left( \pauli{\jrow} \pauli{\irow} \pauli{\jcol} \pauli{\icol} \right)^\dagger } & \tr{\bf A}^* = \tr{{\bf A}^\dagger} \\
&= \frac{1}{4} M_\irow^\icol \, \tr{ \pauli{\icol}  \pauli{\jcol} \pauli{\irow} \pauli{\jrow}  } & {\bf AB}^\dagger = {\bf B}^\dagger {\bf A}^\dagger \\
&= \frac{1}{4} M_\irow^\icol \, \tr{ \pauli{\jcol} \pauli{\irow} \pauli{\jrow} \pauli{\icol} } & \tr{\bf AB} = \tr{\bf BA} \\
&= N_\jcol^\jrow  & \peqn{Mueller_to_target_coherency}
\end{align*}
\end{proof}

Just as the coherency matrix formed by the outer product of a single Jones vector
$\tilde{\cohM{\rho}} = \hat{\Jcol{e}} \otimes \hat{\Jcol{e}}^\dagger$ corresponds to a purely polarized state
\eqnp[see]{spectral_decomposition_of_coherency_matrix}, the target coherency matrix formed by the outer product of a single target vector $\MM{N} = \Sv{k} \otimes \Sv{k}^\dagger / 2$ corresponds to a pure Mueller matrix.
%
Therefore, a Mueller matrix is pure if and only if 
its target coherency matrix, $\MM{N} = \MM{\JT{T}} \dcbo \MM{M}$ 
satisfies \eqn{pure_target_coherency}.
%
\cite{clo86} expressed this definition with the equivalent constraint that a pure Mueller matrix is associated with a target coherency matrix that has only one non-zero eigenvalue, $\lambda$.
%
Combined with its associated eigenvector, $\hat{\Sv{k}}$,
the target vector $\Sv{k} = \lambda^{1/2} \hat{\Sv{k}}$ can be used to compute (up to arbitrary phase) the Jones matrix associated with a pure Mueller matrix,
%
\begin{equation}
\JM{J} = \frac{1}{2} k_\mu \pauli{\mu}.
\end{equation}

\subsection{Impure Mueller Matrices}
\label{app:impure}

%%%%%%%%%%%%%%%%%%%%%%%%%%%%%%%%%%%%%%%%%%%%%%%%%%%%%%%%%%%%%%%%%%%%%%%%%%%%%%
%%%%%%%%%%%%%%%%%%%%%%%%%%%%%%%%%%%%%%%%%%%%%%%%%%%%%%%%%%%%%%%%%%%%%%%%%%%%%%

Some polarimetric transformations are described by impure Mueller matrices that cannot be represented by a linear transformation of the electric field, as described by a Jones matrix.
For example, temporal or spectral depolarization occurs when the response of a system varies as a function of time or frequency on characteristic scales that are smaller than the duration or bandwidth over which the Stokes parameters are integrated.  In this case, the measured Stokes parameters,
%
\begin{equation}
\Sv{S}'=\langle \MM{M}(t,\nu) \rangle \Sv{S},
\end{equation}
where the angle brackets represent integration over time and/or frequency.  

For example, consider stochastic Faraday rotation that causes fluctuations in the position angle $\Delta\Psi$ that are normally distributed with zero mean and standard deviation $\sigma_\Psi$.
The average Mueller matrix that describes the integrated Stokes parameters at a single radio frequency,

\begin{equation} \begin{split}
\langle \MM{R}_{\hat{\Pvs v}}(2\Delta\Psi) \rangle 
&= \left( \begin{array}{cccc}
1 & 0 & 0 & 0 \\
0 & \langle\cos2\Delta\Psi\rangle & -\langle\sin2\Delta\Psi\rangle & 0\\
0 & \langle\sin2\Delta\Psi\rangle & \langle\cos2\Delta\Psi\rangle & 0 \\
0 & 0 & 0 & 1 
\end{array} \right) \\
&= \left( \begin{array}{cccc}
1 & 0 & 0 & 0 \\
0 & d & 0 & 0\\
0 & 0 & d & 0 \\
0 & 0 & 0 & 1 
\end{array} \right)
\end{split} \end{equation}
where the angle brackets represent integration over time and $d=\exp(-2\sigma_\Psi^2) < 1$ describes the depolarization of Stokes Q and U. Note that this Mueller matrix is able to depolarize a purely polarized state, something that cannot be done by a linear transformation of the electric field (\Proof{only_valid}).

\iffalse

The equivalent linear transformation of the coherency matrix is given by
%
\begin{equation}
\dc{ \langle\spinorBilinear{\JM{J}(t,f)}{\JM{J}^\dagger(t,f)}\rangle }{\cohM{\rho}}
\end{equation}
%
where the $\tilde\otimes$ operator represents a tensor product followed by a
transpose over covariant tensor indeces; e.g.,  for matrices \JM{A} and \JM{B},
\begin{equation}
\left\{\spinorBilinear\JM{A}\JM{B}\right\}_{ik}^{jl}
 \equiv 
A_i^l B_k^j. \label{eqn:indexSpinorBilinear} 
\end{equation}
%

\begin{proof}
The tensor product \spinorBilinear\JM{A}\JM{B} has the
following transformation property
%
\begin{align*}
\left( \spinorBilinear\JM{A}\JM{B} \right) \dcbo\,\JM{C}
&= 
\JM{A}\JM{C}\JM{B}. & \text{\citep{car91}}
\end{align*}
% The right-hand side of \eqn{spinorBilinear} is the matrix product of \JM{A}, \JM{B}, and \JM{C}.
Therefore, 
\begin{align*}
    {\JM{J}{\cohM{\rho}}\,\JM{J}}^\dagger
    = \dc{ \left(\spinorBilinear{\JM{J}}{\JM{J}^\dagger}\right) }{\cohM{\rho}}
\end{align*}
%
and the associated Mueller matrix
%
\begin{align*}
  M_\irow^\icol = \frac{1}{2} \dc{\pauli{\irow}}{\dc{ \spinorBilinear{\JM{J}}{\JM{J}^\dagger} }{\pauli{\icol}}}.
\end{align*}
\end{proof}

\fi 
